% !PRINTED 343
\weberpage{343}

Wir wollen hier also nur auf die Tripelgleichungen $9^{\text{ten}}$ Grades
eingehen, und zunächst zeigen, wie sich aus der Tripeleigenschaft
die Bezeichnung und Anordnung der Wurzeln ableiten lässt, die
wir im vorigen Paragraphen für die Wendepunkte kennen gelernt
haben.

Jede der neun Wurzeln kommt in vier Tripeln vor, und im
Ganzen existiren daher $4.9:3=12$ verschiedene Tripel. Nehmen
wir ein beliebiges von diesen Tripeln und bezeichnen dessen
Wurzeln mit
\[
(0,\,0)\ (1,\,2)\ (2,\,1).
\]
Dann sei $(1,\,1)$ eine beliebige vierte Wurzel. Diese bestimmt
mit den drei ersten zunächst drei Tripel, die wir mit
\[
(1,\,1)\ (0,\,0)\ (2,\,2), \quad (1,\,1)\ (1,\,2)\ (1,\,0), \quad (1,\,1)\ (2,\,1)\ (0,\,1)
\]
bezeichnen, und es bleiben noch zwei Wurzeln übrig, die das
vierte Tripel mit $(1,\,1)$ bilden, und das wir nun mit
\[
(1,\,1)\ (2,\,0)\ (0,\,2)
\]
bezeichnen.

Wir suchen nun das durch $(0,\,0), (2,\,0)$ bestimmte Tripel.
In diesem können keine Wurzeln vorkommen, die mit einer der
beiden Wurzeln $(0,\,0), (2,\,0)$ in einer der schon bestimmten Tripel % [sic] page prints feminine "einer" with neuter "Tripel"; expected "einem"
vorkommen, also nicht
\[
(1,\,2)\ (2,\,1)\ (1,\,1)\ (2,\,2)\ (0,\,2),
\]
und es bleibt für die dritte Wurzel dieses Tripels nur $(0,\,1)$
oder $(1,\,0)$ übrig, und dieselben beiden Wurzeln bleiben auch nur
übrig für das durch $(0,\,0)\,(0,\,2)$ bestimmte Tripel. Da wir aber
noch $(0,\,2)$ mit $(2,\,0)$ vertauschen können, so beschränken wir
die Allgemeinheit nicht, wenn wir die zwei weiteren Tripel
\[
(0,\,0)\ (2,\,0)\ (1,\,0), \quad (0,\,0)\ (0,\,2)\ (0,\,1)
\]
annehmen. Nun bestimmen wir die Tripel, die $(1,\,0)$ enthalten,
von denen zwei schon bekannt sind. Die beiden anderen
müssen die Wurzeln $(2,\,2), (0,\,1), (0,\,2), (2,\,1)$ enthalten. Da aber
$(0,\,1)\,(1,\,0)\,(0,\,2)$ und $(0,\,1)\,(1,\,0)\,(2,\,1)$ keine Tripel sind [weil
$(0,\,1)\,(0,\,2)$ das Tripel $(0,\,0)\,(0,\,2)\,(0,\,1)$ und $(0,\,1)\,(2,\,1)$ das
Tripel $(1,\,1)\,(2,\,1)\,(0,\,1)$ bestimmt], so haben wir die weiteren
Tripel
\[
(2,\,2)\ (0,\,1)\ (1,\,0), \quad (1,\,0)\ (0,\,2)\ (2,\,1).
\]
Ebenso ergiebt sich
\[
(0,\,1)\ (2,\,0)\ (1,\,2).
\]
